Paper~7 found that lag-1 rolling-history online recalibration of the
HygienePrio scorer recovers the offline-peek calibration ceiling at
moderate capacity (per-cell recovery ratio $\rho \in \{1.04, 0.99\}$
at $K \in \{50, 100\}$, $\lambda = 3$) but harms performance at
$K = 200$ ($\rho = -0.66$), and named multi-window-history smoothing
as the natural candidate fix. This paper evaluates that fix in a
pre-registered 2{,}250-row five-strategy comparison: fixed (Paper~4
weights), lag1 (Paper~7 baseline), trail3 (equal-weight 3-window
trailing mean), ewma3 (exponentially-weighted 3-window mean,
$\eta = 0.6$), and offline (Paper~7 ceiling). All four pre-registered
hypotheses are rejected, and the direction of rejection is the
contribution: multi-window smoothing does not reverse the
high-capacity hazard --- it makes the hazard substantially worse. At
$K = 100$, EWMA-3's cell-mean recovery ratio is $-1.37$ (versus
lag-1's $+0.99$); at $K = 200$ it is $-0.94$ (versus lag-1's
$-0.66$); at $K = 50$, the matched-capacity moderate regime, EWMA-3
also underperforms lag-1 ($\rho = 0.53$ versus $1.04$). The mechanism
is the inverse of the naive smoothing-helps intuition: at high
capacity the fleet's distribution at window $w$ has shifted
substantially from the three preceding windows, so weighting all
three in the calibration target pulls the selected weights further
from the appropriate-for-window-$w$ optimum. The hazard is monotone
in history length --- less history is better when distributional
change is fast. Deployable online recalibration in this setting
should therefore use the shortest possible history (lag-1 at most)
and be turned off entirely at high capacity. All results are bounded
to the synthetic EEHDA evaluation context.
